Domain randomisation (DR) trains policies on dynamics sampled from a range, but the width of that range is usually chosen by hand. We study, at a deliberately small CPU scale, whether widening the range eventually costs nominal performance faster than it buys robustness. We train a 4-8-1 MLP with the cross-entropy method on a batched copy of gymnasium CartPole in which cart mass, pole mass, pole length and force magnitude are scaled by $\exp(U(-w,w))$, and evaluate on nominal dynamics and on shifts up to $w=2$, which lie outside most training ranges. Over 5 seeds, widening the range monotonically raised mean return under shifted dynamics up to $w\approx1.5$ (from 386.6 with no randomisation to 493.3), and nominal return stayed at the 500-step cap for every width up to $w=3$ (499.9 at $w=3$). We therefore found no width at which nominal performance dropped faster than robustness improved; at most robustness plateaus and slips slightly at $w=3$, within seed noise. A success-gated curriculum that widens the range from zero matched wide fixed DR (470.8 vs 470.4, paired $p=0.96$) and an ungated linear schedule did at least as well, so the gate gave no measurable benefit. The result is limited to one easy task, one learner and 5 seeds and should not be extrapolated to locomotion.
